\section*{Capítulo \ref{ch:g6:earth-a-planet} --- A Terra, um planeta do Sistema Solar}

\begin{solution}{exo:g6:earth-a-planet:1}
\underline{Mercúrio}, \underline{Vênus}, \underline{a Terra},
\underline{Marte}, Júpiter, Saturno, Urano, Netuno. Os \omterm{def:g6:earth-a-planet:planet-types}{planetas rochosos}
são os quatro primeiros.
\end{solution}

\begin{solution}{exo:g6:earth-a-planet:2}
``Um \omterm{def:g6:earth-a-planet:planet-types}{planeta rochoso} é pequeno e denso, com uma superfície sólida. Um
\omterm{def:g6:earth-a-planet:planet-types}{planeta gigante} é muito grande, de baixa densidade, feito sobretudo de
gases e de líquidos.''
\end{solution}

\begin{solution}{exo:g6:earth-a-planet:3}
Vênus (12\,104~km, contra 12\,756~km da Terra). O maior é Júpiter
(142\,984~km); o menor é Mercúrio (4\,879~km).
\end{solution}

\begin{solution}{exo:g6:earth-a-planet:4}
(a) Falso: Vênus não tem água líquida; ele é quente demais. (b) Verdadeiro
($-65$~$^\circ$C em média, contra 15~$^\circ$C). (c) Falso: os quatro
planetas mais próximos do Sol são os \omterm{def:g6:earth-a-planet:planet-types}{planetas rochosos}; os \omterm{def:g6:earth-a-planet:planet-types}{planetas
gigantes} são os quatro mais distantes. (d) Verdadeiro: cerca de
0.7~grama por centímetro cúbico.
\end{solution}

\begin{solution}{exo:g6:earth-a-planet:5}
$4879 \div 12756 \approx 0.382$, ou seja, 38~\%: Mercúrio tem um pouco
mais de um terço da largura da Terra.
\end{solution}

\begin{solution}{exo:g6:earth-a-planet:6}
$120536 \div 12756 \approx 9.4$: cerca de 9 Terras. Saturno é um \omterm{def:g6:earth-a-planet:planet-types}{planeta gigante}.
\end{solution}

\begin{solution}{exo:g6:earth-a-planet:7}
O planeta~A é um \omterm{def:g6:earth-a-planet:planet-types}{planeta gigante}: ele é muito grande e a sua densidade,
1.1~grama por centímetro cúbico, é baixa, como a dos \omterm{def:g6:earth-a-planet:planet-types}{planetas gigantes} do
Sistema Solar (entre 0.7 e 1.7). O planeta~B é um \omterm{def:g6:earth-a-planet:planet-types}{planeta rochoso}: ele é
pequeno e a sua densidade, 5.0, é alta, como a dos \omterm{def:g6:earth-a-planet:planet-types}{planetas rochosos} (entre
3.9 e 5.5).
\end{solution}

\begin{solution}{exo:g6:earth-a-planet:8}
(a) $464 - 15 = 449$~graus mais quente. (b) $15 - (-65) = 80$~graus mais frio.
\end{solution}

\begin{solution}{exo:g6:earth-a-planet:9}
Vênus tem uma atmosfera muito espessa de gás carbônico, que prende o calor
do Sol como um cobertor pesado e não o deixa escapar. Mercúrio, mais perto
do Sol, quase não tem ar, então não consegue guardar o seu calor.
\end{solution}

\begin{solution}{exo:g6:earth-a-planet:10}
$228 - 150 = 78$~milhões de quilômetros mais longe. $78 \div 150 = 0.52$:
a distância de Marte é 52~\% maior que a da Terra.
\end{solution}

\begin{solution}{exo:g6:earth-a-planet:11}
Há bilhões de anos, Marte tinha \omterm{def:g1:landscapes-around-us:water}{rios}, \omterm{def:g1:landscapes-around-us:water}{lagos} e \omterm{def:g5:water-shapes-the-land:flood}{inundações}: ele era mais
quente e mais úmido. Hoje a sua água está congelada, no gelo das calotas
polares e embaixo do chão perto dos \omterm{def:g3:continents-and-oceans:poles}{polos}. A água líquida não consegue
durar na sua superfície agora porque Marte é muito frio ($-65$~$^\circ$C em
média) e porque o seu ar é tão ralo que a água logo congela ou vira vapor.
\end{solution}

\begin{solution}{exo:g6:earth-a-planet:12}
(a) $8.9 \div 9.8 \approx 0.908$, ou seja, 91~\%. (b) Como um gêmeo: quase
o mesmo tamanho (95~\% do diâmetro da Terra), quase a mesma atração da
gravidade, e ele também é um \omterm{def:g6:earth-a-planet:planet-types}{planeta rochoso}. Não como um gêmeo: ele é muito
mais quente (464~$^\circ$C contra 15~$^\circ$C), o seu ar é muito mais
espesso e feito de gás carbônico, com \omterm{def:g4:clouds-rain-snow:cloud}{nuvens} de ácido, e ele não tem água
líquida.
\end{solution}

\begin{solution}{pb:g6:earth-a-planet:1}
\textbf{1.} Vênus (108~milhões~de~km), a Terra (150), Marte (228).

\textbf{2.} $150 - 108 = 42$: Vênus fica 42~milhões de quilômetros mais
perto. $228 - 150 = 78$: Marte fica 78~milhões de quilômetros mais longe.

\textbf{3.} $6792 \div 12756 \approx 0.532$, ou seja, 53~\%: Marte tem
cerca da metade da largura da Terra.

\textbf{4.} Os três são \omterm{def:g6:earth-a-planet:planet-types}{planetas rochosos}: as suas densidades (5.2, 5.5 e
3.9 gramas por centímetro cúbico) são altas, e os seus diâmetros são
pequenos comparados com os dos \omterm{def:g6:earth-a-planet:planet-types}{planetas gigantes}.

\textbf{5.} O ar mais espesso fica em Vênus (92~bar); o mais ralo, em
Marte (0.01~bar).

\textbf{6.} $1 \div 0.01 = 100$: o ar de Marte pressiona cem vezes menos
que o ar da Terra.

\textbf{7.} O seu ar espesso de gás carbônico prende o calor do Sol e não
o deixa escapar, como um cobertor pesado.

\textbf{8.} $92 \div 1 = 92$: no chão de Vênus, o ar pressiona
\textbf{92 vezes} mais forte que no chão da Terra. Uma sonda precisa ter
paredes muito resistentes para não ser esmagada (e também precisa resistir
ao calor e às \omterm{def:g4:clouds-rain-snow:cloud}{nuvens} de ácido).

\textbf{9.} A temperatura média da Terra, 15~$^\circ$C, fica entre
0~$^\circ$C e 100~$^\circ$C: a água não está nem congelada nem fervendo,
então ela continua líquida.

\textbf{10.} Em Vênus, a 464~$^\circ$C, muito acima de 100~$^\circ$C, a
água ferveria na mesma hora e viraria vapor. Em Marte, a $-65$~$^\circ$C,
muito abaixo de 0~$^\circ$C, ela congelaria (e viraria vapor aos poucos,
porque o ar é muito ralo).

\textbf{11.} Há muito tempo, Marte era mais quente e mais úmido que hoje:
a água líquida podia correr na sua superfície. O seu \omterm{def:g5:weather-and-climate:climate}{clima} mudou muito:
hoje ele é um \omterm{def:g5:weather-and-climate:desert}{deserto} frio e seco.

\textbf{12.} Por exemplo: ``Queridas sondas, aqui na Terra temos água
líquida, em \omterm{def:g2:water-on-earth:ocean}{oceanos}, \omterm{def:g1:landscapes-around-us:water}{rios} e \omterm{def:g1:landscapes-around-us:water}{lagos}. As nossas temperaturas são amenas,
15~$^\circ$C em média, entre o congelamento e a fervura da água. Estamos na
distância certa do Sol, nem perto demais nem longe demais. E a nossa
atmosfera nos mantém aquecidos, dá aos seres vivos os gases de que eles
precisam e deixa a água continuar líquida. É por isso que a vida é
possível aqui.''
\end{solution}
